\chapter{Published Baseline: Approximate GOOP Solving}

\section{Chapter role}

This chapter documents the completed and published foundation of the thesis:
\emph{Approximate solutions to games of ordered preference}
\cite{delasheras:2025}. The rest of the manuscript extends this baseline rather
than replacing it.

\section{Core idea}

The published method combines receding horizon with lexicographic IBR over time
for approximate GOOP solving. The motivation is to avoid the cost of repeatedly
solving a tightly coupled multi-agent lexicographic program at each control
cycle.

\section{Algorithmic structure}

For each receding-horizon step:

\begin{enumerate}
    \item initialize trajectory predictions for all agents from shifted previous
    solutions,
    \item run bounded IBR rounds,
    \item in each round, each agent solves its lexicographic best-response
    subproblem against fixed predictions,
    \item aggregate resulting controls and execute the next segment.
\end{enumerate}

This process keeps preference semantics explicit while reducing per-step
computational burden.

\section{Implementation and tooling}

The implementation is in Julia \cite{julia} and uses the GOOP ecosystem around
TrajectoryGamesBase.jl \cite{peters:2024a}, ParametricMCPs.jl
\cite{peters:2025}, and PATH-backed complementarity solving
\cite{dirkse:1995}. The released package is LEXIBROverTime.jl \cite{code}.

\section{Representative results}

The paper reports substantial runtime gains compared to a tightly coupled
baseline while preserving approximate solution quality over practical iteration
budgets.

\begin{table}[tbp]
\centering
\caption{Representative values reported in \cite{delasheras:2025}.}
\label{tab:published_runtime}
\begin{tabular}{@{}lrrrr@{}}
\toprule
 & \multicolumn{2}{c}{$K=2$} & \multicolumn{2}{c}{$K=3$} \\
\cmidrule(r){2-3}\cmidrule(l){4-5}
Method & Solve time [s] & $L_1$ distance & Solve time [s] & $L_1$ distance \\
\midrule
Baseline & 44.51 & -- & 2676.00 & -- \\
$L=1$ & 0.36 & $3.28\times10^{-4}$ & 0.32 & $5.45\times10^{-4}$ \\
$L=2$ & 0.59 & $1.22\times10^{-5}$ & 0.63 & $6.56\times10^{-5}$ \\
$L=5$ & 1.39 & $1.78\times10^{-6}$ & 1.20 & $4.96\times10^{-5}$ \\
\bottomrule
\end{tabular}
\end{table}

\section{What is reused in this thesis}

This baseline contributes the reusable components adopted in later chapters:

\begin{itemize}
    \item receding-horizon execution interface,
    \item trajectory warm-start logic,
    \item lexicographic objective structure,
    \item solver hyperparameterization and diagnostics.
\end{itemize}

\section{Open points left by the paper}

The paper leaves two practical gaps addressed by the MSc thesis:

\begin{enumerate}
    \item deployment-level integration with physical/simulated robots,
    \item inverse estimation of hidden preference structure.
\end{enumerate}

These are exactly the focus of the following chapters.
